\section{Three-slot packing for Iwaniec's remainder coefficients}\label{sec:pack}

This section proves a general packing theorem. It then verifies the
hypotheses for the remainder coefficients of Iwaniec's linear sieve
\cite[Theorem~4]{AA37}. The packing and coefficient identities provide
the modulus factorization used in the prime-discrepancy estimates below.
The functions below specialize Lichtman's optimization
\cite[Proposition~6.1, (6.2), (6.4)]{Lichtman23} to spectral exponent
$7/32$ and moving-residue exponent $1$. We prove the labelled-atom
statement here; its two-bin step is related to the corrected argument
in Pascadi's Lemma~5.6 \cite{Pascadi25}.

\subsection{A combinatorial packing theorem}

Put
\[
 \overline H=\frac{153}{256},\qquad \mu=\frac{25}{128},
 \qquad \vartheta=\frac7{32}.
\]
Here $\vartheta$ is the fixed spectral exponent $2\theta_{\rm RP}$.
Later in this section $\theta$ denotes the sieve level.
For $x\ge0$, define
\[
 v(x)=\min\left\{\frac{1+x}{2},\frac{32-34x}{39},
                         \frac{57}{64}-\frac32x\right\}.
\]
For real $x$ and $y\ge0$, put $\psi(x,y)=x\mathbf 1_{x\ge y}$.
Also put
\[
 w(a,b,c)=\psi\left(\min\left\{
       \frac{132-68c}{214},\frac{153-96c}{256},1-2b\right\},
       \frac{1+a}{2}\right).
\]
For an ordered triple $a\ge b\ge c\ge0$, set
\[
\begin{split}
 H(a,b,c)=\max\{&v(a),v(b),v(a+b),v(a+b+c),w(a,b,c),w(a,c,b),\\
                 &\psi(v(a+c),a+2b+c),
                   \psi(v(b+c),2a+b+c)\},
\end{split}
\]
Define $\widetilde H(a,b,c)=H(a,b,c)$ if $a\le\mu$.
If $a>\mu$, put $\widetilde H(a,b,c)=v(a)$. For a finite non-increasing
list $e_1\ge\cdots\ge e_m\ge0$, let $\bar e_i=e_i$ for $i\le m$
and $\bar e_i=0$ for $i>m$. Repeated values are distinct labelled
atoms. The empty list is allowed.

A list $f_1\ge\cdots\ge f_m\ge0$ satisfies $\mathrm{WF}(y)$ if
\[
 f_1+\cdots+f_{j-1}+2f_j\le y\qquad(1\le j\le m).
 \tag{4.1}\label{eq:pack-wf}
\]
Write
\[
 b_-=2h-1,\qquad b_A=\frac{32-39h}{34},\qquad
 b_B=\frac{2-\vartheta-2h}{3},\qquad b_+=\min(b_A,b_B).
\]

\refstepcounter{theorem}\textbf{Theorem~\thetheorem\ (three-slot packing).}\label{thm:pack-a}
Let $0<\delta\le3\cdot10^{-9}$,
$1/2\le h\le\overline H$, $d_*=h-2\delta$, and
$\delta\le g\le1/3-\delta/2$. Suppose the labelled list
$e_1\ge\cdots\ge e_m\ge0$ satisfies
\[
 \mathrm{WF}(d_*)\quad\hbox{and}\quad
 h\le\widetilde H(\bar e_1,\bar e_2,\bar e_3).
 \tag{4.2}\label{eq:pack-hypotheses}
\]
Then the deterministic rule $\mathcal R$ below partitions the labels
into $J_1\sqcup J_2\sqcup J_3$. If
$\Sigma_s=\sum_{i\in J_s}e_i$, then
\[
\begin{aligned}
 \Sigma_1&\le g,\\
 2g+\Sigma_2+2\Sigma_3&\le1-4\delta,\\
 \vartheta+g+\Sigma_1+5\Sigma_2+2\Sigma_3&\le2-4\delta,\\
 (1+2\vartheta)g+\Sigma_1+(5-2\vartheta)\Sigma_2
       +(2+2\vartheta)\Sigma_3&\le2-4(1+\vartheta)\delta.
\end{aligned}
 \tag{4.3}\label{eq:pack-loads}
\]
The output depends only on the labelled list and on $h,\delta,g$.

\refstepcounter{theorem}\label{lem:pack-1}\textbf{Lemma~\thetheorem.} For $x\ge0$,
$v(x)=\min\{(1+x)/2,57/64-3x/2\}$. Thus $v$ is concave,
$v\le\overline H$, and $\widetilde H\le\overline H$.

\emph{Proof.} Consider the middle affine function in the definition
of $v$. Its difference with $(1+x)/2$ is $(25-107x)/78$, which is
nonnegative for $x\le25/107$. Its difference with $57/64-3x/2$ is
$(1568x-175)/2496$, nonnegative for $x\ge175/1568$.
As $175/1568<25/107$, it is never the strict minimum of all three.
The minimum of the remaining two affine functions is concave and has
maximum $(1+\mu)/2=\overline H$. Each $w$ is zero or at most
$(153-96c)/256\le\overline H$. Every $\psi$ branch is zero or a value
of $v$. This proves both upper bounds. \hfill$\square$

\refstepcounter{theorem}\label{lem:pack-2}\textbf{Lemma~\thetheorem.} For $1/2\le h\le\overline H$ and
$0<\delta\le3\cdot10^{-9}$:
\begin{enumerate}
\item $h\le v(b)$ for $b\ge0$ exactly when $b_-\le b\le b_+$.
\item $b_+=b_B$, and $b_A>b_B$.
\item $b_+-b_-=(153-256h)/96\ge0$ and
      $b_+-\mu=(153/128-2h)/3\ge0$.
\item $d_*-b_+-2b_-\ge3/256-2\delta>0$ and
      $d_*-3b_-\ge3/256-2\delta>0$.
\item $\vartheta+2(d_*+2\delta)+3b_B=2$ and
      $(2+2\vartheta)(d_*+2\delta)+(3-4\vartheta)b_A=2$.
\end{enumerate}

\emph{Proof.} The three affine branches of $v$ give
$b\ge2h-1$, $b\le b_A$, and $b\le b_B$.
These are closed inequalities. The difference
$b_A-b_B=(567/16-49h)/102$ decreases with $h$ and at
$h=\overline H$ is $525/8704>0$, proving (2). The formulas in
(3) and the first bound in (4) follow by expansion. Their minima occur
at $h=\overline H$. For the second bound in (4), use
$d_*-b_-=1-h-2\delta$ and $3-5\overline H=3/256$.
For (5), use $2+2\vartheta=39/16$,
$3-4\vartheta=17/8$, and
$(3-4\vartheta)b_A=2-39h/16$. Finally
$3/256>6\cdot10^{-9}\ge2\delta$. \hfill$\square$

\refstepcounter{theorem}\label{lem:pack-3}\textbf{Lemma~\thetheorem.} If $h>0$ and $w(x,y,z)\ge h$, then
$z\le b_+-b_-$.

\emph{Proof.} Since $h>0$, the activation in $\psi$ is on. Hence
$(132-68z)/214\ge h$ and $(153-96z)/256\ge h$. These give
$z\le b_A-b_-$ and $z\le b_B-b_-$, respectively. \hfill$\square$

\refstepcounter{theorem}\label{lem:pack-4}\textbf{Lemma~\thetheorem\ (two bins).} Let $c_1,c_3\ge0$ and initial loads
$u_0\le c_1$, $w_0\le c_3$. Suppose a finite list
$f_1,\ldots,f_k\ge0$ satisfies
\[
 u_0+w_0+\sum_{i<j}f_i+2f_j\le c_1+c_3\qquad(1\le j\le k).
 \tag{4.4}\label{eq:pack-two-bin}
\]
Process the atoms in order. Place each in bin 1 if it fits, and
otherwise in bin 3 if it fits. Every atom is placed, and the final loads
are at most $c_1,c_3$.

\emph{Proof.} Suppose $f_j$ is the first atom fitting neither bin.
Then the current loads $U,W$ satisfy
$U+W=u_0+w_0+\sum_{i<j}f_i$, $U+f_j>c_1$, and
$W+f_j>c_3$. Adding the last two inequalities contradicts
\eqref{eq:pack-two-bin}. \hfill$\square$

The rule permits zero atoms and repeated values. Its three cases are:
\begin{enumerate}
\item If $g+b_+\ge d_*$, apply Lemma~\ref{lem:pack-4} to the whole list.
      Use bin caps $g,d_*-g$ for slots 1 and 2, and leave slot 3 empty.
\item If $g+b_+<d_*$, test in order
      $K=\{1\},\{1,2\},\{1,2,3\},\{1,3\},\{2,3\}$. Put
      $K'=K\cap\{1,\ldots,m\}$ and $b=\sum_{i\in K'}e_i$.
      Call $K$ legal if $b_-\le b\le b_+$ and the remaining atoms,
      in their original order, satisfy $\mathrm{WF}(d_*-b)$. For
      the first legal $K$, place $K'$ in slot 2. Apply Lemma~\ref{lem:pack-4}
      to the rest in slots 1 and 3, with caps $g,d_*-g-b$.
\item If no $K$ is legal, use caps $g,b_-,d_*-g-b_-$ for slots
      1, 2, 3. Put atom 1 in slot 2 if it exists. Process subsequent
      atoms in order, placing each in the first slot in which it fits.
      Put the first atom fitting no slot in slot 2. Call the
      resulting slot-2 load $b$. Apply Lemma~\ref{lem:pack-4} to the remaining
      atoms in slots 1 and 3, with caps $g,d_*-g-b$.
\end{enumerate}

\textbf{Proof of Theorem~\ref{thm:pack-a}.} Write $\mathcal R$ for the rule just
defined. Since $h\ge1/2$, $b_-\ge0$; also
$d_*=h-2\delta\ge1/2-6\cdot10^{-9}>1/3\ge g$.
It is enough to prove one of the following envelopes:
\[
\begin{split}
 (E_{\rm glob}):&\quad g\ge d_*-b_+,\quad
 \Sigma_1\le g,\quad\Sigma_2\le d_*-g,\quad\Sigma_3=0;\\
 (E_{\rm crit}(b)):&\quad b_-\le b\le b_+,\quad
 \Sigma_1\le g,\quad\Sigma_2\le b,\quad
 \Sigma_3\le d_*-g-b.
\end{split}
\]
The second, third, and fourth left sides in \eqref{eq:pack-loads}
have nonnegative load coefficients. Under $E_{\rm crit}(b)$,
Lemma~\ref{lem:pack-2}(5) gives
\[
\begin{aligned}
2g+\Sigma_2+2\Sigma_3&\le2d_*-b\le2d_*-b_-=1-4\delta,\\
\vartheta+g+\Sigma_1+5\Sigma_2+2\Sigma_3
 &\le\vartheta+2d_*+3b
 \le\vartheta+2d_*+3b_B=2-4\delta.
\end{aligned}
\]
For the fourth load, similarly,
\[
\begin{aligned}
 &(1+2\vartheta)g+\Sigma_1+(5-2\vartheta)\Sigma_2
       +(2+2\vartheta)\Sigma_3\\
 &\qquad\le(2+2\vartheta)d_*+(3-4\vartheta)b\\
 &\qquad\le(2+2\vartheta)d_*+(3-4\vartheta)b_A\\
 &\qquad=2-4(1+\vartheta)\delta.
\end{aligned}
\]
Here $3-4\vartheta=17/8>0$. Under $E_{\rm glob}$, use
$g\ge d_*-b_+$ in the two expressions whose coefficient of $g$
is negative. This gives
\[
\begin{aligned}
2g+\Sigma_2&\le g+d_*
 \le\tfrac13-\tfrac\delta2+\overline H-2\delta
 =\tfrac{715}{768}-\tfrac52\delta\le1-4\delta,\\
\vartheta+2g+5\Sigma_2&\le\vartheta+5d_*-3g
 \le\vartheta+2d_*+3b_+=2-4\delta,\\
(2+2\vartheta)g+(5-2\vartheta)\Sigma_2
 &\le(5-2\vartheta)d_*-(3-4\vartheta)g\\
 &\le(2+2\vartheta)d_*+(3-4\vartheta)b_+
 \le2-4(1+\vartheta)\delta.
\end{aligned}
\]
The first line uses $1-715/768=53/768>3\delta/2$.
In both envelopes $\Sigma_1\le g$, proving the first inequality
of \eqref{eq:pack-loads}. Equality at the case boundary belongs to
case 1; all displayed endpoint inequalities are weak.

Consider case 1. The clause $\mathrm{WF}(d_*)$ of
\eqref{eq:pack-hypotheses} is exactly the hypothesis
\eqref{eq:pack-two-bin} for the whole list. Here $u_0=w_0=0$ and the
total cap is $d_*$.
Lemma~\ref{lem:pack-4} gives $E_{\rm glob}$. Now suppose
$g+b_+<d_*$. If a tested set $K$ is legal, the remaining list
satisfies \eqref{eq:pack-two-bin} with total cap $d_*-b$. Lemma~\ref{lem:pack-4}
then gives $E_{\rm crit}(b)$.

We verify the rule when no tested set is legal. Two deletion facts
will be used. First let $K'$ be a leading set $\{1,\ldots,k\}$. Subtract
$b=\sum_{i\le k}e_i$ from every later inequality in
$\mathrm{WF}(d_*)$. This proves $\mathrm{WF}(d_*-b)$ for the remaining
list. Next let $K'$ be $\{1,3\}$, $\{2,3\}$, or, when $m=2$, the
truncation $\{2\}$ of $K=\{2,3\}$. Let $b_-\le b\le b_+$.
The only unselected atom among the first three has index $i_0=2$ or $1$.
If it is below $b_-$, it is first in the remaining list, and
\[
 2e_{i_0}+b<2b_-+b_+\le d_*
\]
by Lemma~\ref{lem:pack-2}(4). When $m=2$ this is the whole remaining list.
Now take a remaining rank at least four. Its original
prefix sum contains both selected atoms. Subtracting $b$ proves the
reduced mask. Thus the nonconsecutive set is legal whenever these
conditions hold.

Assume now that no tested set is legal. First, $\bar e_1\le\mu$. Otherwise the definition of $\widetilde H$ and
\eqref{eq:pack-hypotheses} give $h\le v(e_1)$. Then
$b_-\le e_1\le b_+$ by Lemma~\ref{lem:pack-2}(1), and $\{1\}$ is legal.
The same lemma and $b_+\ge\mu$ show $\bar e_1<b_-$; otherwise
$\{1\}$ is legal. Hence every padded coordinate is below $b_-$.
Since $\bar e_1\le\mu$, the branch condition gives
$\widetilde H=H$. A branch in its maximum is at least $h$.
Neither $v(\bar e_1)$ nor $v(\bar e_2)$ can be that branch.
Each would put the corresponding coordinate at least $b_-$. A branch
$v(\bar e_1+\bar e_2)$ or
$v(\bar e_1+\bar e_2+\bar e_3)$ would make a leading set legal.
If
\[
 \psi(v(\bar e_1+\bar e_3),\bar e_1+2\bar e_2+\bar e_3)\ge h,
\]
then $v(\bar e_1+\bar e_3)\ge h$, and
$b=\bar e_1+\bar e_3\in[b_-,b_+]$. The only unselected atom among
the first three is $e_2<b_-$. So the deletion fact makes $\{1,3\}$
legal, a contradiction. The other $\psi$ branch similarly makes
$\{2,3\}$ legal. Therefore either
$w(\bar e_1,\bar e_2,\bar e_3)\ge h$ or
$w(\bar e_1,\bar e_3,\bar e_2)\ge h$. Lemma~\ref{lem:pack-3} gives, respectively,
\[
 \bar e_3\le b_+-b_-,\qquad\hbox{or}\qquad
 \bar e_2\le b_+-b_-.
 \tag{4.5}\label{eq:pack-w-tail}
\]

In case 3, the slot-3 cap is
\[
d_*-g-b_-=1-h-g-2\delta\ge
1-\overline H-\frac13-\frac32\delta
=\frac{53}{768}-\frac32\delta>0.
\]
Thus all three caps are nonnegative. Atom 1 fits slot 2 because
$e_1<b_-$. In the first $w$ case, atom 2 cannot
be a stopping atom. Slots 1 and 3 are empty. If $e_2$ fit neither,
then $2e_2>d_*-b_-\ge2b_-$ by Lemma~\ref{lem:pack-2}(4), contrary to
$e_2<b_-$. Thus any stopping atom has index $j\ge3$ and
$e_j\le e_3\le b_+-b_-$. In the second $w$ case, a stopping atom
has $j\ge2$ and $e_j\le e_2\le b_+-b_-$. If there is no stop,
the final loads are bounded by $(g,b_-,d_*-g-b_-)$. This is
$E_{\rm crit}(b_-)$, also when $m=0$.

If the stop is at $j$, let $(u_0,v_0,w_0)$ be the loads before
$e_j$. After placing $e_j$ in slot 2, its load $b=v_0+e_j$ satisfies
\[
 b_-<b\le b_-+(b_+-b_-)=b_+.
\]
The lower bound holds because $e_j$ did not fit slot 2. The upper
bound follows from \eqref{eq:pack-w-tail}. The $j$-th inequality in
$\mathrm{WF}(d_*)$ gives
$u_0+w_0+e_j\le d_*-b$. Since $e_j$ did not fit slot 1,
$u_0+e_j>g$, so $w_0<d_*-g-b$. For every later rank, the original
mask is exactly the two-bin condition with total cap $d_*-b$.
Lemma~\ref{lem:pack-4} completes the assignment and gives $E_{\rm crit}(b)$.
The rule processes every atom once and has fixed tie priorities, so it
is well defined. All cases imply \eqref{eq:pack-loads}. \hfill$\square$

\subsection{A local ceiling and a small-atom criterion}

In this subsection write $\theta_{\rm RP}$ for the exponent in the
unramified Ramanujan bound and $\vartheta_{\rm RP}=2\theta_{\rm RP}$. The
Kim--Sarnak bound \cite{KS03} gives $\theta_{\rm RP}=7/64$ and
$\vartheta_{\rm RP}=7/32=\vartheta$.

\refstepcounter{theorem}\textbf{Proposition~\thetheorem\ (local ceiling and small-atom packing).}\label{prop:pack-ceiling}
Suppose
$0\le\vartheta_{\rm RP}\le7/32$, that is $0\le\theta_{\rm RP}\le7/64$,
and $0\le t\le1/3$. Let
$x,y,z_3\ge0$ satisfy the zero-loss local slot constraints
\[
\begin{aligned}
 x-t&\le0, &2t+y+2z_3&\le1,\\
 \vartheta_{\rm RP}+t+x+5y+2z_3&\le2,
 &(1+2\vartheta_{\rm RP})t+x+(5-2\vartheta_{\rm RP})y+(2+2\vartheta_{\rm RP})z_3&\le2,
\end{aligned}
\]
and $x+y+z_3\le1-t$. Then their maximum total load is
\[
 H_{\vartheta_{\rm RP}}=\frac{5-\vartheta_{\rm RP}}{8}=\frac{5-2\theta_{\rm RP}}8.
\]
For fixed $t$, equality holds precisely at
$x=t$, $y=(1-\vartheta_{\rm RP})/4$, $z_3=(3+\vartheta_{\rm RP})/8-t$.

Suppose instead that $0\le e_j\le u$ with $u\ge0$. Let the
non-increasing labelled atoms obey $\mathrm{WF}(h-2\delta)$, with
\[
 \frac12\le h\le H_{\rm small}(\vartheta_{\rm RP},u):=\frac{5-\vartheta_{\rm RP}-3u}{8},
 \qquad 0<\delta\le\frac1{100},\qquad 0\le g\le\frac13.
\]
Then they admit a three-slot partition with the four load
inequalities \eqref{eq:pack-loads}. Here $\vartheta_{\rm RP}$ replaces $\vartheta$.
In particular, $H_{\rm small}(7/32,1/500)=19101/32000$ and
$H_{7/32}=153/256$.

\emph{Proof.} Seven times $x-t\le0$, three times
$2t+y+2z_3\le1$, and once $\vartheta_{\rm RP}+t+x+5y+2z_3\le2$ give
$8(x+y+z_3)+\vartheta_{\rm RP}\le5$. Equality forces equality in these three
constraints, giving the stated point. Conversely, that point has
$y\ge25/128$, $z_3\ge1/24$, and total load plus $t$ at most
$23/24$. Substitution in the fourth constraint gives
$2-3\vartheta_{\rm RP}(1-\vartheta_{\rm RP})/4\le2$, proving the maximum and its equality cases.

For the small-atom assertion put
\[
\begin{aligned}
 d&=h-2\delta,& b_-&=2h-1,&
 b_B&=\frac{2-\vartheta_{\rm RP}-2h}{3},\\
 b_A&=\frac{2-(2+2\vartheta_{\rm RP})h}{3-4\vartheta_{\rm RP}},&
 b_+&=\min(b_A,b_B).
\end{aligned}
\]
The denominators are positive, and direct subtraction gives
\[
 b_B-b_-=\frac{5-\vartheta_{\rm RP}-8h}{3},\qquad
 b_A-b_-=\frac{5-4\vartheta_{\rm RP}-(8-6\vartheta_{\rm RP})h}{3-4\vartheta_{\rm RP}}.
\]
Hence $u\le b_+-b_-$ is equivalent to
\[
 h\le\frac{5-\vartheta_{\rm RP}-3u}{8},\qquad
 h\le\frac{5-4\vartheta_{\rm RP}-(3-4\vartheta_{\rm RP})u}{8-6\vartheta_{\rm RP}}.
\]
The second upper bound exceeds the first by
$[6\vartheta_{\rm RP}(1-\vartheta_{\rm RP})+14\vartheta_{\rm RP} u]/[8(8-6\vartheta_{\rm RP})]\ge0$.
So the assumed bound on $h$ gives $u\le b_+-b_-$. Start with capacities
$(g,b_-,d-g-b_-)$. They are nonnegative, since the last is at least
$1-5/8-1/3-2\delta\ge1/24-1/50>0$. Place atoms in order in the
first slot in which they fit. If all fit, take $b=b_-$. Otherwise, put the first failed atom $e_j$ in slot 2. Let $b$ be the
new slot-2 load. Its failure there and $e_j\le u$ give
$b_-<b\le b_-+u\le b_+$. Let $(U,V,W)$ be the previous loads.
Then $\mathrm{WF}(d)$ gives $U+W+e_j\le d-b$. Failure in slot 1 gives
$W<d-g-b$. For each later atom, subtract the fixed slot-2 load $b$
from its original $\mathrm{WF}(d)$ inequality. Lemma~\ref{lem:pack-4} places the
remaining atoms in slots 1 and 3 with capacities $g,d-g-b$.
Both cases reach
\[
 \Sigma_1\le g,\qquad\Sigma_2\le b,\qquad
 \Sigma_3\le d-g-b,\qquad b_-\le b\le b_+.
\]
The last three left sides of \eqref{eq:pack-loads}, with $\vartheta_{\rm RP}$ in
place of $\vartheta$, are at most
\[
 \begin{gathered}
 2d-b_-=1-4\delta,\qquad
 \vartheta_{\rm RP}+2d+3b_B=2-4\delta,\\
 (2+2\vartheta_{\rm RP})d+(3-4\vartheta_{\rm RP})b_A=2-4(1+\vartheta_{\rm RP})\delta,
 \end{gathered}
\]
respectively. This proves the partition claim. \hfill$\square$

These are a relaxed local spectral ceiling and a sufficient uniform
small-atom packing criterion. The actual local inequalities have
positive losses, so the relaxed ceiling is a supremum. For example,
take the losses $29\sigma$, $8\sigma$, and $15\sigma$ in the first
three slot constraints. They give
\[
 x+y+z_3\le H_{\vartheta_{\rm RP}}-\frac{242\sigma}{8}
                    -\frac{\vartheta_{\rm RP}\sigma}{160}
\]
when the moving-frequency factor is $X^{\vartheta_{\rm RP}(1+\sigma/20)}$.
The small-atom bound also needs strict room in an actual source.
If $h=\rho \theta+2\delta\le H_{\rm small}(\vartheta_{\rm RP},u)$ and
$\delta>0$, then $\theta<H_{\rm small}(\vartheta_{\rm RP},u)$.
The next subsection applies this criterion to the structured
coefficient family of Iwaniec's construction.

\subsection{Endpoint geometry for Iwaniec's coefficients}

Fix the parameters in Theorem~\ref{thm:aa37-interior} and put
\[
\begin{gathered}
 \tau=\sigma^9,\qquad \rho=1+\tau,\qquad
 \Theta(u)=H_{\rm small}(\vartheta,u)=\frac{5-\vartheta-3u}{8},\\
 \eta_\Theta=\Theta(u)-\theta,\\
 \delta=30\sigma,\qquad d_*=\rho \theta,\qquad h=d_*+2\delta.
\end{gathered}
 \tag{4.6}\label{eq:pack-aa37-parameters}
\]
Proposition~\ref{prop:pack-ceiling} gives the sufficient packing threshold $\Theta(u)$.
The relaxed local ceiling is $H_\vartheta=153/256$. The fixed
support exponent for these coefficients is $\Theta(u)$, and
$\Theta(u)<H_\vartheta$ because $u>0$. The margin
$\eta_\Theta=\Theta(u)-\theta$ leaves room for the smooth divisor and tail-prime
ranges. Since the theorem also imposes $D^{\sigma^2}\le z$, every
nonempty tail has the required roughness floor. Lemma~\ref{lem:aa37-pack} verifies
$\rho \theta+60\sigma<\Theta(u)$ and the hypotheses of
Theorem~\ref{thm:pack-a}.
Let $D=N^{\theta}$, $z=N^u$, and $Y=D^{\sigma^2}$. Write
$\mathcal P_N=\{p\text{ prime}:p\nmid2N\}$ and
$P_N(Y)=\prod_{p<Y,\ p\in\mathcal P_N}p$. Fix an Iwaniec
sign. Write $\mathcal D^\nu(D)$ for its finite family of nonempty
label vectors $\mathbf D=(D_1,\ldots,D_k)$. They have
$D_k\le\cdots\le D_1<D^{1/2}$, grid coordinates in
$\{D^{\sigma^2\rho^j}:j\ge0\}$, and the cubic conditions
\[
\begin{array}{ll}
 \nu=+:&D_1\cdots D_{2j}D_{2j+1}^3<D
                    \quad(2j+1\le k),\\
 \nu=-:&D_1\cdots D_{2j-1}D_{2j}^3<D
                    \quad(2j\le k).
\end{array}
 \tag{4.7}\label{eq:pack-aa37-cubic}
\]
The prime labels in a vector range over
$D_j\le p_j<\min(D_j^\rho,z)$, with $p_j\nmid2N$. A vector with
an empty coordinate interval contributes zero and may be omitted.

\setcounter{lemma}{\value{theorem}}
\refstepcounter{lemma}\textbf{Lemma~\thelemma\ (Iwaniec coefficient packing).}\label{lem:aa37-pack}
For every contributing vector and
every $\xi\in[60\sigma,1/3+14\sigma]$, let
$e_j=\log_N\min(D_j^\rho,z)$, sorted in non-increasing order.
Then $\mathrm{WF}(d_*)$ holds and
\[
 h\le\widetilde H(\bar e_1,\bar e_2,\bar e_3).
 \tag{4.8}\label{eq:pack-aa37-hyp}
\]
Consequently Theorem~\ref{thm:pack-a} partitions the tail labels into three slots.
Assign the Iwaniec smooth divisor $d_{AA}<D^\sigma$ to slot 1.
Let $Q,S,R$ be the resulting lower dyadic slot scales. Then
\[
\begin{aligned}
 \log_N Q&\le \xi-\left(1+\frac{7271}{256}\right)\sigma,\\
 \log_N(T^2SR^2)&\le1-60\sigma,\\
 \log_N(A_{\max}^{\vartheta}TQS^5R^2)
   &\le2-\left(15+\frac{95221}{1280}\right)\sigma,\\
 \log_N\left((TR/S)^{2\vartheta}TQS^5R^2\right)
   &\le2-\left(15+\frac{22407}{256}\right)\sigma,
\end{aligned}
\tag{4.9}\label{eq:pack-scales}
\]
where $T=N^\xi$ and $A_{\max}=X^{1+\sigma/20}$ for $X\le N$.

\emph{Proof.} Put $s_j=\log_ND_j$. The grid lower bound gives
$s_j\ge\sigma^2\theta$. The cubic conditions in
\eqref{eq:pack-aa37-cubic} imply $\sum_{i<j}s_i+2s_j<\theta$
for every rank $j$. At a cubic rank this follows by dropping one
nonnegative copy of $s_j$. At an intermediate rank,
\[
 \sum_{i<j}s_i+2s_j
 =\sum_{i<j-1}s_i+s_{j-1}+2s_j
 \le\sum_{i<j-1}s_i+3s_{j-1}<\theta,
\]
using $s_{j-1}\ge s_j$ and the cubic condition at rank $j-1$.
For the first rank in the minus case use $2s_1<\theta$, which follows
from $D_1<D^{1/2}$. Thus $\sum_{i<j}e_i+2e_j\le
\rho(\sum_{i<j}s_i+2s_j)<\rho \theta=d_*$.
Also $e_j\le u$, so $\bar e_1\le u<\mu$.
For this range, the branch $w(\bar e_1,\bar e_2,\bar e_3)$ is active.
Its activation threshold is at most $(1+u)/2\le0.501$, while
\[
 \min\left\{\frac{132-68u}{214},
              \frac{153-96u}{256},1-2u\right\}
 =\frac{153-96u}{256}
 =H_{\rm small}(\vartheta,u)=\Theta(u).
\]
Each branch in the minimum decreases with its argument, and
$\bar e_2,\bar e_3\le u$. Take $0\le u\le1/500$. Each other affine branch minus
$(153-96u)/256$ is nonnegative at both endpoints. Hence it is
nonnegative on the whole interval. In particular, the
minimum defining $w(\bar e_1,\bar e_2,\bar e_3)$ is at least
$\Theta(u)>0.501$, so its activation condition holds. Thus
$\widetilde H\ge\Theta(u)$. Next,
$\tau \theta+60\sigma<\eta_\Theta$, because
$60\sigma\le0.3\eta_\Theta$ and
$\tau \theta<0.6\sigma^9\le0.6\sigma<0.7\eta_\Theta$.
Therefore
\[
 \frac12<h=\rho \theta+60\sigma
 =\Theta(u)-\eta_\Theta+\tau \theta+60\sigma
 <\Theta(u)<\overline H.
\]
Also $0<\delta\le3\cdot10^{-9}$. The range for $\xi$ gives
$g=\xi-\delta\in[\delta,1/3-\delta/2]$.
Theorem~\ref{thm:pack-a} applies.

The smooth-only term is the empty-list case. Here
$\widetilde H(0,0,0)=H(0,0,0)=\overline H$, since the $w$ branch
equals $\min\{132/214,153/256,1\}=\overline H$ and its activation
threshold is $1/2$. Thus Theorem~\ref{thm:pack-a} applies to the empty list as well.
The smooth divisor remains in slot 1, and slots 2 and 3 are units.

For the scale bounds, the slot-1 smooth divisor contributes at most
$\sigma \theta<\sigma\overline H$ to the logarithmic load. The other
slot loads are bounded by $\Sigma_2,\Sigma_3$. Thus
\[
\begin{aligned}
 \log_NQ&\le\Sigma_1+\sigma\overline H,&
 \log_NS&\le\Sigma_2,&
 \log_NR&\le\Sigma_3,\\
 \log_NT&=\xi=g+\delta,&
 \log_NA_{\max}&\le1+\sigma/20.
\end{aligned}
\]
The first line of \eqref{eq:pack-loads} gives
\[
 \log_NQ-\xi\le-(30-\overline H)\sigma
 =-\left(1+\frac{7271}{256}\right)\sigma.
\]
The second gives
\[
 \log_N(T^2SR^2)\le1-2\delta=1-60\sigma.
\]
For the third scale, \eqref{eq:pack-loads} gives
\[
 \log_N(A_{\max}^{\vartheta}TQS^5R^2)
 \le2-3\delta+\left(\frac{\vartheta}{20}+\overline H\right)\sigma
 =2-\left(15+\frac{95221}{1280}\right)\sigma.
\]
For the fourth, it gives
\[
 \log_N((TR/S)^{2\vartheta}TQS^5R^2)
 \le2-(3+2\vartheta)\delta+\overline H\sigma
 =2-\left(15+\frac{22407}{256}\right)\sigma.
\]
The exact constants in \eqref{eq:pack-scales} hold for the lower scales.
We round each of the three aggregate slot scales $Q,R,S$ once, to a
dyadic representative within a factor two. Each monomial below has
total degree at most $8$ in $Q,R,S$. So rounding multiplies it by at
most $2^8$. Individual prime labels are not rounded.
This adds at most
$8\log 2/\log N$ to its base-$N$ exponent. Fix $\sigma>0$. For sufficiently large $N$, each positive margin
absorbs this cost. The tightest margin, for $Q<TX^{-29\sigma}$, is
$103\sigma/256-8\log 2/\log N>0$. \hfill$\square$

For fixed $B>0$ and $N/(2(\log N)^B)\le X\le N$, put
$\omega=\log_N(N/X)$. Then
$\omega\le(B\log\log N+\log2)/\log N$, so eventually
$\log_XN=(1-\omega)^{-1}<1+0.01\sigma$. The first exponent in
\eqref{eq:pack-scales}, after dividing by $T$, is negative. So
conversion from base $N$ to base $X$ improves that relative inequality.
Each other exponent is at most $2$;
conversion costs at most $0.02\sigma$. The displayed margins absorb
this cost.

For every contributing vector, the last-rank mask gives
$\sum_js_j<\theta$. Since each $s_j\ge\sigma^2\theta$, it also gives
$k<\sigma^{-2}$. Thus $\prod_jD_j<D$, and $p_j<D_j^\rho$ gives
\[
 d\prod_jp_j<D^\sigma D^\rho=D^{1+\sigma+\tau}
 =N^{(1+\sigma+\tau)\theta}<N^{\Theta(u)},
\]
where $(\sigma+\tau)\theta<\eta_\Theta$.
All prime factors avoid $2N$. Thus the full coefficient is supported
on $q<N^{\Theta(u)}$ with $(q,2N)=1$. Its coefficients are
$O_\sigma(1)$. The $Y$-smooth part of a supported modulus determines
$d$. The remaining labelled prime assignments number at most $k!$,
and $k<\sigma^{-2}$. For each sign the Iwaniec family has fewer than
$e^{8\sigma^{-3}}$ vectors. So the sum over vectors keeps a bound
depending only on $\sigma$. The smooth-only term satisfies the same
support and coefficient bounds.

The cubic-to-quadratic implication just proved is uniform in the vector
length. We record the two bounds obtained above:
\[
 k<\sigma^{-2},\qquad \sum_{j=1}^k s_j<\theta.
 \tag{4.10}\label{eq:pack-aa37-length}
\]
These bounds also control the coefficient regrouping below.

\subsection{Three-slot coefficient identity}

Fix a sign $\nu$ and a vector $\mathbf D$. The set of
equal-label pairs in the Iwaniec grid is
$I(\mathbf D)=\{(i,j):1\le i<j\le k,\ D_i=D_j\}$. Distinct grid labels have
disjoint half-open prime intervals
$[D_i,\min(D_i^\rho,z))$. Adjacent grid values differ by
the power $\rho$. Thus repeated primes can occur only in
$I(\mathbf D)$. Define
\[
 K_{\mathbf D}(\mathbf p)=
   \prod_{(i,j)\in I(\mathbf D)}(1-\mathbf1_{p_i=p_j}).
\]
For the Iwaniec coefficients $\lambda_d^\nu(\mathbf D)$, define the
unrestricted coefficient
\[
 C_{\mathbf D}^{\rm full}(q)=
 \sum_{\substack{d<D^\sigma\\d\mid P_N(Y)}}
 \lambda_d^\nu(\mathbf D)
 \sum_{\substack{D_j\le p_j<\min(D_j^\rho,z)\\p_j\nmid2N}}
 \mathbf1_{q=d\prod_{j=1}^k p_j}.
 \tag{4.11}\label{eq:pack-aa37-vector}
\]
The masked coefficient $C_{\mathbf D}^{\rm mask}$ is defined by
inserting $K_{\mathbf D}(\mathbf p)$ in the inner sum. The coefficient
of the smooth-only term, the sole empty-tail term in the Iwaniec formula,
is
$C_{\emptyset}(q)=\sum_{d<D^\sigma,d\mid P_N(Y)}
\phi_d^\nu\mathbf1_{q=d}$. Thus the full signed coefficient is
\[
 w_N^\nu(q)=C_{\emptyset}(q)+
       \sum_{\mathbf D\in\mathcal D^\nu(D)}
             C_{\mathbf D}^{\rm mask}(q).
 \tag{4.12}\label{eq:pack-aa37-weight}
\]

\refstepcounter{lemma}\textbf{Lemma~\thelemma\ (slot factorization and collision support).}\label{lem:aa37-slot}
Fix $\xi$ in
the range of Lemma~\ref{lem:aa37-pack} and let $J_1\sqcup J_2\sqcup J_3$ be its
packing partition. All $d$ and $p_j$ sums below use the ranges in
\eqref{eq:pack-aa37-vector}. Define
\[
\begin{aligned}
 \gamma_1(m)&=\frac1{|J_1|!}
    \sum_{d,(p_j)_{j\in J_1}}
       \lambda_d^\nu(\mathbf D)\mathbf1_{m=d\prod_{j\in J_1}p_j},\\
 \gamma_s(m)&=\frac1{|J_s|!}
    \sum_{(p_j)_{j\in J_s}}
       \mathbf1_{m=\prod_{j\in J_s}p_j},\qquad s=2,3.
\end{aligned}
\]
For $s=2,3$ set $\gamma_s=\mathbf1_{m=1}$ when $J_s$ is empty;
$\gamma_1$ always retains the smooth-divisor sum. Then
\[
 C_{\mathbf D}^{\rm full}(q)=C_{\rm slot}
       (\gamma_1*\gamma_2*\gamma_3)(q),\qquad
 C_{\rm slot}=|J_1|!|J_2|!|J_3|!,
 \tag{4.13}\label{eq:pack-aa37-slot-id}
\]
and $|\gamma_s(m)|\le1$. Consider the full Iwaniec family. The difference between the masked
and unrestricted coefficients has pointwise size $O_\sigma(1)$. It is
supported on moduli divisible by $p^2$ for some prime $p\ge Y$.

\emph{Proof.} A labelled tuple $(d,(p_j)_{j\le k})$ splits uniquely
into its three slot tuples. Multiplying the three normalized finite
sum definitions by $C_{\rm slot}$ therefore gives
\eqref{eq:pack-aa37-slot-id}. This regroups a finite sum and
allows repeated primes. Fix a product $m$. Its prime factors
below $Y$ determine $d$ uniquely: $d\mid P_N(Y)$, and every tail
prime is at least $D_j\ge Y$. Once $d$ is fixed, the labelled prime
assignments within a slot number at most $|J_s|!$. The Iwaniec bounds
$|\lambda_d^\nu(\mathbf D)|\le1$ give $|\gamma_s(m)|\le1$;
the same argument applies to $\gamma_1$. The smooth-only term has
the corresponding bound from $|\phi_d^\nu|\le1$. Since
$k<\sigma^{-2}$, the scalar
$C_{\rm slot}=|J_1|!|J_2|!|J_3|!$ is $O_\sigma(1)$.

Suppose the distinctness mask removes a tuple. Then some selected pair
has $p_i=p_j\ge Y$, and its modulus is divisible by $p_i^2$. For fixed
$\sigma$, \eqref{eq:pack-aa37-length} bounds the number of labels.
The Iwaniec family has fewer than $e^{8\sigma^{-3}}$ vectors for each
sign. The number of labelled
assignments to a fixed modulus is at most $k!$. Hence the aggregate
coefficient difference is $O_\sigma(1)$.

The prime-source cost of this collision support is estimated in
(B.5)--(B.6). \hfill$\square$

The smooth-only term $C_{\emptyset}$ is already a one-slot coefficient.
Take $J_1=J_2=J_3=\emptyset$, put
$\gamma_1(m)=\sum_{d<D^\sigma,d\mid P_N(Y)}
\phi_d^\nu\mathbf1_{m=d}$, and
$\gamma_2=\gamma_3=\mathbf1_{m=1}$. Its coefficient identity and
pointwise bound follow directly. The full coefficient $w_N^\nu$ in
\eqref{eq:pack-aa37-weight} is the sum of this term and the vector
terms above. Every partition is fixed
before the prime realizations are summed.

This section proves the packing inequalities and the slot identity for
the unrestricted coefficient. It also gives the collision support
needed in Appendix~B.
The prime-source estimate is proved there.
